\documentclass[letterpaper,10pt]{article}

% =========================================================================
% CONFIGURAZIONE FONT E CODIFICA (Ottimizzata per ATS e font moderni)
% =========================================================================
\usepackage[T1]{fontenc}
\usepackage[utf8]{inputenc}    % Consente l'uso diretto di accenti (es. à, è, ì)

\usepackage{latexsym}
\usepackage[empty]{fullpage}
\usepackage{titlesec}
\usepackage{marvosym}
\usepackage[dvipsnames]{xcolor} % Sostituito 'color' con il più moderno 'xcolor'
\usepackage{verbatim}
\usepackage{enumitem}
\usepackage[hidelinks]{hyperref}
\usepackage{fancyhdr}
\usepackage[english]{babel}
\usepackage{tabularx}
\input{glyphtounicode}

% Forza la generazione di un PDF leggibile dagli ATS (Text-searchable)
\pdfgentounicode=1

% --- Opzione A: ROBOTO (Attuale) ---
%~ \usepackage[sfdefault]{roboto}

% --- Opzione B: INTER (Super moderno/Tech) ---
\usepackage[sfdefault]{inter}

% --- Opzione C: HELVETICA (Classico Corporate) ---
%~ \usepackage[scaled=0.92]{helvet}
%~ \renewcommand{\familydefault}{\sfdefault}

% --- Opzione D: Lato (Molto elegante, leggermente più morbido e umano) ---
%~ \usepackage[default]{lato}

% --- Opzione E: EB GARAMOND (Elegante/Accademico con le grazie) ---
%~ \usepackage{ebgaramond}

% --- Opzione F: Charter (Solido/Robusto e leggibile anche a dimensioni ridotte) ---
%~ \usepackage{charter}

% --- Opzione G: TeX Gyre Pagella (Una bellissima reinterpretazione del Palatino, molto equilibrato) ---
%~ \usepackage{tgpagella}



\addtolength{\oddsidemargin}{-0.5in}
\addtolength{\evensidemargin}{-0.5in}
\addtolength{\textwidth}{1.0in}
\addtolength{\topmargin}{-0.5in}
\addtolength{\textheight}{1.0in}

\urlstyle{same}

% =========================================================================
% STILE DELLE SEZIONI (Personalizzato NavyBlue)
% =========================================================================
\definecolor{NavyBlue}{RGB}{0, 32, 96}
\titleformat{\section}{
  \vspace{-5pt}\scshape\raggedright\large\bfseries\color{NavyBlue}
}{}{0em}{}[\color{NavyBlue}\titlerule \vspace{-4pt}]

% =========================================================================
% STILE DELLE SEZIONI (Non personalizzato NavyBlue)
% =========================================================================
%~ \titleformat{\section}{
  %~ \vspace{-5pt}\scshape\raggedright\large\bfseries
%~ }{}{0em}{}[\color{black}\titlerule \vspace{-4pt}]

% =========================================================================
% MACRO PERSONALIZZATE (Ambiente di sviluppo del layout)
% =========================================================================
\newcommand{\resumeItem}[1]{
  \item\small{
    {#1 \vspace{-2pt}}
  }
}

\newcommand{\resumeSubheading}[4]{
  \vspace{-1pt}\item
    \begin{tabular*}{0.97\textwidth}[t]{l@{\extracolsep{\fill}}r}
      \textbf{#1} & #2 \\
      \textit{\small#3} & \textit{\small #4} \\
    \end{tabular*}\vspace{-5pt}
}

\newcommand{\resumeProjectHeading}[2]{
    \item
    \begin{tabular*}{0.97\textwidth}{l@{\extracolsep{\fill}}r}
      \small#1 & #2 \\
    \end{tabular*}\vspace{-5pt}
}

\renewcommand{\labelitemii}{$\circ$}

% =========================================================================
% INIZIO DEL CORPO DEL CV
% =========================================================================
\begin{document}

% --- HEADER ---
\begin{center}
    \textbf{\Huge \scshape Leonardo Pozzati} \\ \vspace{4pt}
    \small Tilburg, Netherlands $|$ +39 334 710 8340 $|$ +31 6 29808909 \\ \vspace{2pt}
    \href{mailto:leon.pozzati@gmail.com}{leon.pozzati@gmail.com} $|$ 
    \href{https://github.com/vaeVict1s}{github.com/vaeVict1s} $|$ 
    \href{https://www.linkedin.com/in/leonardo-pozzati-4769471bb}{linkedin.com/in/leonardo-pozzati}
\end{center}

% --- PROFESSIONAL SUMMARY ---
\section*{Profile}
\small{Senior Software Engineer and Systems Engineer with over 10 years 
of experience specializing in low-latency, performance-critical systems 
across fintech and high-tech industries. Proven track record of building 
proprietary FX trading infrastructure processing 0.8M+ messages/sec 
with sub-250$\mu$s latency. 
\newline
Combines a rigorous academic foundation in Mathematical Physics with 
deep expertise in C++17/20, lock-free concurrency, and Python data analysis pipelines.}

% --- EDUCATION ---
\section{Education}
  \begin{itemize}[leftmargin=0.15in, label={}]
    \resumeSubheading
      {Università degli Studi Roma Tre}{Roma, Italy}
      {Master's Degree in Mathematical Physics -- 110/110 cum Laude}{}
      \begin{itemize}[leftmargin=0.1in]
        \resumeItem{Research on the relaxation to the equilibrium of 
        supersaturated gases through critical droplets in metastable thermodynamic systems.}
        \resumeItem{Designed and implemented stochastic and deterministic 
        mathematical models; developed high-performance C++ simulations 
        and executed massive data analysis pipelines in Python within 
        Non-Equilibrium Statistical Mechanics.}
      \end{itemize}
    \resumeSubheading
      {Università degli Studi Roma Tre}{Roma, Italy}
      {Bachelor's Degree in Physics}{}
  \end{itemize}

% --- TECHNICAL SKILLS ---
\section{Technical Skills}
 \begin{itemize}[leftmargin=0.15in, label={}]
    \small{\item{
     \textbf{Programming Languages:}{ C, C++, Python, Bash, SQL, Linux, OOP and UML.} \\
     \textbf{Quantitative \& Analytics:}{ Mathematical Physics, Statistical Mechanics, 
     Time-Series Analysis, Machine Learning (MLP, LSTM, TCN, Transformers), Signal Evaluation.} \\
     \textbf{Low-Latency \& Systems:}{ Multi-threading, Lock-free Concurrency, 
     Memory Management (NUMA-aware, custom cache-aligned allocators), 
     Linux Kernel Tuning, Linux \texttt{epoll} event loops.} \\
     \textbf{Tools \& Infrastructure:}{ ZeroMQ, CI/CD (Jenkins), Git, 
     Docker, Profiling (\texttt{perf}, Flamegraphs, \texttt{heaptrack}), Agile Methodologies.}
    }}
 \end{itemize}

% --- EXPERIENCE ---
\section{Professional Experience}
  \begin{itemize}[leftmargin=0.15in, label={}]

    \resumeSubheading
      {Thermo Fisher Scientific}{Eindhoven, Netherlands}
      {\textbf{Software Designer}}{2025 -- Present}
      \begin{itemize}
        \resumeItem{Develop and maintain core automation software for T
        ransmission Electron Microscopes (TEMs) applied to semiconductor 
        chip anomaly detection.}
        \resumeItem{Collaborate on low-level architectural design, 
        development, and testing of efficiency-critical software modules 
        utilizing C/C++, Python, and Bash in Agile environments.}
      \end{itemize}

    \resumeSubheading
      {ASML}{Veldhoven, Netherlands}
      {\textbf{Senior Software Engineer} - Universal Light System}{2024}
      \begin{itemize}
        \resumeItem{Joined a core team responsible for building a new 
        universal light system for ASML’s lithography machines.}
        \resumeItem{Designed and implemented software modules using C/C++, 
        with scripting in Python and Bash.}
        \resumeItem{Followed Agile practices to coordinate development, 
        testing, and integration from the ground up.}
      \end{itemize}
      
    \resumeSubheading
      {ASML}{Veldhoven, Netherlands}
      {\textbf{Senior Software Engineer} - Safety and Diagnostic}{2022 - 2024}
      \begin{itemize}
        \resumeItem{Assigned to the core Safety and Diagnostic team for 
        the ASML's High-NA EUV lithography platform.}
        \resumeItem{Improved system robustness by integrating, refactoring, 
        and testing both new and legacy code.}
        \resumeItem{Collaborated with the T\&I (Test \& Integration) team 
        to integrate CI/CD pipelines and streamline deployments in the 
        way of working of the team.}
        \resumeItem{Authored the software design documentation for 
        safety-critical components.}
        \resumeItem{Collaborated with mechanical engineers in the testing 
        and validation of multiple physical machine subsystems, 
        providing software support, diagnostics and issue investigation.}
      \end{itemize}
    
    \newpage
    \resumeSubheading
      {ProMarket}{Rome, Italy}
      {\textbf{Senior Software Engineer / Quant Developer}}{2018 -- 2021}
      \begin{itemize}
        \resumeItem{Built and owned proprietary high-frequency FX trading 
        infrastructure supporting internal prop trading teams deploying 
        firm capital.}
        \resumeItem{Designed and implemented a low-latency routing and 
        messaging framework using C++17, Python, and ZeroMQ 
        (pub/sub, dealer/router, binary star), achieving consistent 
        production latencies of 150--250$\mu$s.}
        \resumeItem{Reduced end-to-end median messaging latency by 50\% 
        (from 400$\mu$s to 200$\mu$s) via CPU pinning, NUMA-aware workload 
        partitioning, \texttt{boost::lockfree::queue}, and lean 
        serialization formats.}
        \resumeItem{Scaled messaging throughput from 400K to over 0.8M 
        messages/sec by implementing epoll-based event loops 
        \newline
        (\texttt{EPOLLONESHOT}), custom cache-aligned allocators, 
        and continuous profiling using \texttt{perf}, Flamegraphs, 
        and \texttt{heaptrack}.}
        \resumeItem{Refactored the real-time risk-checking module with 
        a lock-free memory pool, cutting heap allocation volume under 
        peak load by 30\% and stabilizing memory footprint.}
        \resumeItem{Resolved market-open production latency spikes by 
        replacing inefficient \texttt{gettimeofday()} calls with 
        \newline
        \texttt{clock\_gettime(CLOCK\_MONOTONIC\_RAW)}, successfully 
        reducing latency jitter by 40\%.}
        \resumeItem{Developed orchestration and coordination services 
        for internal strategy deployment, state management, and system 
        diagnostics using Python, Bash, and Docker.}
        \resumeItem{Contributed to a proprietary financial sentiment 
        analysis engine, extracting quantitative trading signals from 
        market data and news feeds to evaluate pre-production signal quality.}
        \resumeItem{Collaborated directly with traders during initial 
        production rollouts, aligning system performance and operational 
        requirements with trading desk needs.}
        \resumeItem{Supported live production for 18+ months, managing 
        10--15K orders/day across 3--5 FX trading venues with multiple 
        parallel execution strategies.}
      \end{itemize}

  \end{itemize}

% --- PERSONAL PROJECTS ---
\section{Key Projects}
  \begin{itemize}[leftmargin=0.15in, label={}]
    \resumeProjectHeading
      {\textbf{WOPR} $|$ \emph{Personal Quant \& ML Trading Project}}{}
      \begin{itemize}
        \resumeItem{Engineered an end-to-end quantitative pipeline in 
        Python for processing CME sessionized Nasdaq futures data.}
        \resumeItem{Implemented range-bar feature engineering, 
        Time-to-Horizon (TTH) label generation, and trained multi-model 
        deep learning architectures including MLP, LSTM, TCN, and Transformers.}
        \resumeItem{Enhanced pipeline reproducibility and system robustness 
        by unifying core labeling logic and integrating checkpoint-driven 
        scaling and metadata tracking across train, evaluation, and 
        inference cycles.}
      \end{itemize}
    \resumeProjectHeading
      {\textbf{Domain-Specific Chatbot Training} $|$ \emph{Educational Software Engineering Notebook}}{}
      \begin{itemize}
        \resumeItem{Developing a structured, reference-heavy knowledge 
        base and data compilation pipeline focused on optimizing 
        LLMs for specific engineering domains.}
      \end{itemize}
  \end{itemize}

% --- LANGUAGES ---
\section{Languages}
 \begin{itemize}[leftmargin=0.15in, label={}]
   \small{\item{
     \textbf{Italian:}{ Native} $|$ 
     \textbf{English:}{ C1 - Academic IELTS} $|$ 
     \textbf{Dutch:}{ Self-studying}
   }}
 \end{itemize}

\end{document}
